% !TEX root = ../../main.tex

\section{Evaluación de Candidatos}

Cada permutación válida se entrega al algoritmo existente de \texttt{ApplicativeDo}. En otras palabras, la extensión no define un nuevo algoritmo de construcción de planes applicative; reutiliza el mecanismo de \textsf{GHC} sobre distintas secuencias de entrada. El resultado de cada evaluación es un árbol de statements junto con su costo.

En el ejemplo guía, la evaluación produce la siguiente comparación:

\begin{center}
\begin{tabular}{clcl}
\hline
Candidato & Orden & Plan conceptual & Costo \\
\hline
0 & \texttt{[1,2,3,4]} & $s_1 ; (s_2 \mid (s_3 ; s_4))$ & 3 \\
1 & \texttt{[1,3,2,4]} & $(s_1 \mid s_3) ; (s_2 \mid s_4)$ & 2 \\
2 & \texttt{[1,3,4,2]} & $(s_1 \mid s_3) ; (s_4 \mid s_2)$ & 2 \\
3 & \texttt{[3,1,2,4]} & $(s_3 \mid s_1) ; (s_2 \mid s_4)$ & 2 \\
4 & \texttt{[3,1,4,2]} & $(s_3 \mid s_1) ; (s_4 \mid s_2)$ & 2 \\
\hline
\end{tabular}
\end{center}

El candidato original alcanza costo 3, que coincide con \texttt{ApplicativeDo} normal sobre el orden escrito por el programador. Los candidatos restantes alcanzan costo 2 porque exponen dos niveles de paralelismo applicative: primero entre los productores independientes y luego entre los consumidores que ya tienen disponibles sus entradas.

Esta evaluación es el punto donde se integra la contribución con el estado del arte. La extensión amplía el conjunto de órdenes posibles; \texttt{ApplicativeDo} sigue siendo el componente que decide cómo agrupar cada orden y cuánto cuesta el plan resultante.
